\section{Cross-Platform Application Design}
\label{sec:application}

GrabDrop is implemented as two native clients---Android and
Desktop---that interoperate seamlessly via a shared network protocol.

\subsection{Android Client}

The Android client is built with Kotlin and Jetpack Compose
(Material~3), targeting Android~10+ (API~29). The architecture incorporates several key components: a foreground service designated \texttt{GrabDropService} that hosts the camera pipeline, gesture detector, network manager, and overlay manager while operating independently of the UI lifecycle; CameraX integration providing the \texttt{ImageAnalysis} use case for front camera frames at configurable frame rates; MediaPipe tasks-vision (v0.10.14) for hand landmark detection on each camera frame, operating in \texttt{VIDEO} mode to enable temporal tracking; ONNX Runtime Android (v1.20.0) for executing the pruned and quantized TCN model during gesture classification; MediaProjection for capturing screenshots upon \texttt{GRAB} events; and window-manager-based overlay animations that provide visual feedback including a flash effect on grab, a ripple effect on release, and a wakeup indicator.

\subsection{Desktop Client}

The Desktop client is a Python CLI application compatible with Linux,
macOS, and Windows. Its architecture mirrors the Android client,
incorporating OpenCV for camera capture via \texttt{VideoCapture}, MediaPipe Python (v0.10.14) for hand landmark detection, ONNX Runtime ($\geq$1.17.0) for TCN inference, platform-specific screenshot tools including \texttt{spectacle} for KDE/Wayland, \texttt{grim} for generic Wayland, \texttt{scrot} for X11, \texttt{screencapture} for macOS, or \texttt{mss} for Windows, and Tkinter-based overlays for visual feedback indicators.

\subsection{Network Protocol}

Both clients share a unified zero-configuration network protocol
operating over the local area network:

\subsubsection{Device Discovery (UDP)}

Each device broadcasts a JSON heartbeat message every 3~seconds via
both multicast (\texttt{239.255.77.88:9877}) and broadcast
(\texttt{255.255.255.255:9877}) for reliability. Messages contain a
unique device ID, device name, and timestamp. Devices not heard from
for 10~seconds are removed from the discovery table.

\subsubsection{Screenshot Transfer (UDP + TCP)}

Upon a \texttt{GRAB} event, the system initiates a three-phase transfer protocol. First, the sender captures a screenshot and starts a TCP server on a random port. Second, a \texttt{SCREENSHOT\_READY} JSON message is broadcast via UDP, containing the TCP port and file size. Third, when a receiver performs a \texttt{RELEASE} gesture, it connects to the sender's TCP server, sends the command \texttt{"GET\textbackslash n"}, and receives a 4-byte big-endian length header followed by the raw PNG data.

\subsubsection{Retroactive Matching}

To handle network latency, the protocol supports retroactive matching:
if a \texttt{RELEASE} gesture occurs before the screenshot offer
arrives, the release timestamp is recorded and matched against
incoming offers within a 3-second window.

\subsection{Cross-Platform Interoperability}

Android and Desktop clients auto-discover each other on the same LAN
and transfer screenshots transparently. The protocol is
platform-agnostic: any device implementing the UDP/TCP wire format
(Section~7 of the protocol specification) can join the GrabDrop
network.
